\clearpage
\fancyhead[L]{\small CERES.}
\hypertarget{native-317-lower}{}
\begin{center}
\textbf{CERES.}
\par\vspace{4pt}\hrule\vspace{4pt}
\fontsize{10}{12}\selectfont Göttingische gelehrte Anzeigen. Stück 87. S. 857..860. 1809. Juni 3.
\par\vspace{4pt}\hrule\vspace{4pt}
\end{center}

\par\indent Im 40. Stück unserer vorigjährigen Anzeigen [S. 299. d. B.] haben wir die zwölften Elemente der \emph{Ceres} mitgetheilt, welche Hr. Prof. \textsc{Gauss} auf die Beobachtungen vom Jahre 1807 gegründet hatte; wir fahren fort, diejenigen Resultate zu geben, welche nach den so eben geschlossenen Rechnungen aus den Beobachtungen des Jahrs 1808 hervorgegangen sind. Der tiefe Stand des Planeten im Steinbocke hat die Beobachtungen für unsere nordlichen Gegenden sehr erschwert; auf der hiesigen Sternwarte konnte nur eine einzige Beobachtung im Meridian angestellt werden, die überdiess nicht für vorzüglich gut gelten kann.


\begin{center}\fontsize{10}{12}\selectfont\setlength{\tabcolsep}{5pt}\renewcommand{\arraystretch}{1.15}\begin{tabular}{lr} 1808 Mittl. Zeit Juli 25. & $13^{\text{u}}\;14^{\text{m}}\;29^{\text{s}}$.\\ Gerade Aufst. & $322^\circ\;8'\;57''6$.\\ Südl. Abweich. & $27^\circ\;52'\;12''5$\\
\end{tabular}\end{center}


\SourceLeaf{318}\par\indent Weiter sind bisher keine Beobachtungen bekannt geworden, als die, welche Hr. v. \textsc{Lindenau} auf der Seeberger Sternwarte, und Hr. \textsc{Santini} in Padua angestellt haben. Mit diesen Hülfsmitteln hat Hr. Prof. \textsc{Gauss} die Opposition von 1808, den Fehler der XII. Elemente und die Correctionen bestimmt, welche diesen Elementen zugefügt werden müssen, um mit allen bisher beobachteten sechs Oppositionen in die möglich genaueste Uebereinstimmung gebracht zu werden. Wir hoffen, dass den Astronomen die Mittheilung dieser Resultate in unsern Blättern bei der jetzt so sehr erschwerten Correspondenz nicht unwillkommen sein wird.

\par\indent Die Opposition wurde nach möglich sorgfältigster Discussion der Beobachtungen gefunden, wie folgt:

\begin{center} 1808 Aug. 5. \quad $11^{\text{u}}\ 30^{\text{m}}\ 13^{\text{s}}$ mittl. Zeit in Göttingen

wahre Länge der Ceres .\;.\;.\;\ $313^\circ\ 13'\ 40''2$

wahre geocentrische Breite .\;.\;\ $12\ 43\ 44{,}8$ südlich
\end{center}

\par\indent Die XII. Elemente weichen diesmal nur Eine Minute in der Länge, und ein paar Secunden in der Breite ab: ein ganz unbedeutender Unterschied, wenn man bedenkt, dass die angewandten Störungen noch nicht vollständig sind, und daher eine vollkommene Vereinigung aller in dem Zeitraume von beinahe acht Jahren gemachten Beobachtungen ohnehin als unmöglich angesehen werden muss. Hr. Prof. \textsc{Gauss}, welcher für die künftige vollständigere Berechnung der Störungen alles Nöthige bereits entworfen hat, aber freilich die Zeit der wirklichen Ausführung dieser eben so delicaten als weitläuftigen Arbeit noch nicht festsetzen kann, hat es inzwischen für interessant gehalten, zu untersuchen, wie genau sich die sämmtlichen sechs beobachteten Oppositionen noch mit seinen im Jahre 1802 berechneten Störungstafeln vereinigen lassen, und der Erfolg zeigt, dass auch diesmal die Differenzen noch so sehr unerheblich sind, dass um ihretwillen die vollständigern Störungen immer noch eine Zeitlang entbehrt werden können. Man darf doch hiebei nicht vergessen, dass künftig, wenn diese Störungen vollständig zugezogen werden müssen, die Berechnung Eines Planetenorts gewiss vier Mal so viele Zeit und Arbeit kosten wird, als jetzt. Die neuen Elemente der Ceres, welche von den XII. nur äusserst wenig abweichen, sind folgende, wo die Epochen für den Meridian von Göttingen gelten:

\Needspace{23\baselineskip}
\begin{center} XIII. \emph{Elemente der Ceres.}
\end{center}


\begin{center}\fontsize{10}{12}\selectfont\setlength{\tabcolsep}{5pt}\renewcommand{\arraystretch}{1.15}\begin{tabular}{crr}  & Mittlere Länge. & Sonnennähe.\\\hline 1801 & $77^\circ\;18'\;36''5$ & $146^\circ\;26'\;0''1$\\ 1802 & $155\;28\;23.4$ & $\phantom{146}\;28\;1.4$\\ 1803 & $233\;38\;10.3$ & $\phantom{146}\;30\;2.6$\\ 1804 & $312\;0\;48.1$ & $\phantom{146}\;32\;4.2$\\ 1805 & $30\;10\;35.0$ & $\phantom{146}\;34\;5.4$\\ 1806 & $108\;20\;21.8$ & $\phantom{146}\;36\;6.6$\\ 1807 & $186\;30\;8.7$ & $\phantom{146}\;38\;7.9$\\ 1808 & $264\;52\;46.5$ & $\phantom{146}\;40\;9.5$\\ 1809 & $343\;2\;33.4$ & $\phantom{146}\;42\;10.7$\\
\end{tabular}\end{center}

\begin{center}\fontsize{9}{11}\selectfont\setlength{\tabcolsep}{5pt}\renewcommand{\arraystretch}{1.15}\begin{tabular}{lr} Tägliche mittlere tropische Bewegung .\; & $770''9230$\\ Excentricität 1806 .\;.\;.\;.\;.\;.\;.\;.\; & $0.0785028$\\ Jährliche Abnahme .\;.\;.\;.\;.\;.\;.\;.\; & $0.00000583$\\ Logarithm der halben grossen Axe .\;.\; & $0.4420486$\\ Aufsteigender Knoten 1806 .\;.\;.\;.\; & $80^\circ\;53'\;41''3$\\ Jährliche Bewegung .\;.\;.\;.\;.\;.\;.\;.\; & $+1{,}48$\\ Neigung der Bahn 1806 .\;.\;,\;.\;.\; & $10^\circ\;37'\;31''2$\\ Jährliche Abnahme .\;.\;.\;.\;.\;.\;.\;.\;.\;.\; & $0{,}44$\\
\end{tabular}\end{center}


\Needspace{13\baselineskip}\par\indent Folgendes ist die Uebereinstimmung dieser neuen Elemente mit allen sechs bisher beobachteten Oppositionen.

\SourceLeaf{319}\begin{center} Unterschied:
\end{center}


\begin{center}\fontsize{10}{12}\selectfont\setlength{\tabcolsep}{5pt}\renewcommand{\arraystretch}{1.15}\begin{tabular}{crr}  & \shortstack{heliocentrische Länge\\in der Bahn} & \shortstack{heliocentrische\\Breite}\\\hline 1802 & $-0''5$ & $-1''9$\\ 1803 & $+22.1$ & $-5.3$\\ 1804 & $-34.9$ & $-1.2$\\ 1806 & $+26.7$ & $-3.6$\\ 1807 & $-26.9$ & $-11.7$\\ 1808 & $+20.8$ & $-6.0$\\
\end{tabular}\end{center}


\par\noindent Es muss hierbei bemerkt werden, dass die Störungen der Länge und des Radius Vectors bei dieser Rechnung aus den Störungstafeln entlehnt sind, welche Hr. Prof. \textsc{Gauss} im Märzhefte der Monatl. Corresp. [1803 März. Seite 235 u. f. dieses Bandes] bekannt gemacht hat; die Breitenstörungen hingegen sich auf eine spätere vollständigere Entwickelung gründen, deren Resultate bisher noch nicht gedruckt sind. Die Geringfügigkeit der Unterschiede bei den Breiten darf man als eine Bestätigung ansehen, dass bei jener Bestimmung keine erhebliche Glieder übersehen sind. [Der Inhalt dieser Mittheilung ist in der Monatlichen Correspondenz Band XIX. Seite 504\,..\,512. 1809 Mai wiedergegeben. Dort sind auch die von \textsc{Gauss} berechneten Ephemeriden, für Ceres 1809 Juli 18 bis 1810 April 16, für Juno von 1809 September 4 bis 1810 Mai 10 abgedruckt. Sie enthalten für Mitternacht in Göttingen mit viertägigen Zwischenräumen geocentrische Gerade Aufsteigung und geocentrische Abweichung mit Minuten, den Abstand von der Erde auf vier Decimalen genau.]

\clearpage
\fancyhead[L]{\small GAUSS AN VON LINDENAU.}
\begin{center}
\textbf{GAUSS AN VON LINDENAU.}
\par\vspace{4pt}\hrule\vspace{4pt}
\fontsize{10}{12}\selectfont Monatliche Correspondenz. Band XX. S. 78, 79. 1809 Juli.
\par\vspace{4pt}\hrule\vspace{4pt}
\end{center}
\begin{flushright} Göttingen, den 14. August 1809.
\end{flushright}

\par\indent ---\enspace---\enspace--- Was ich Ihnen heute mittheilen kann, sind ein paar Beobachtungen der neuen Planeten, welche ich in diesen Tagen gemacht habe. Vorzüglich war es mir darum zu thun, die Pallas bei Zeiten aufzusuchen, weil ich voriges Jahr in Ermangelung guter \emph{vollständiger} Beobachtungen (noch jetzt fehlen mir zuverlässigere \emph{Declinationen}, und ich fürchte fast, dass ich keine anderen als meine eignen und die in Mailand am Aequatorialsector beobachteten werde benutzen können) die Elemente nicht hatte verbessern können, und wir uns also diesmal mit der Ephemeride behelfen müssen, welche Hr. Prof.
\textsc{Bode} im Jahrbuche 1811 nach den im Jahr 1807 bestimmten Elementen berechnet hat. --- Indess hat es eben keine Mühe gekostet den Planeten darnach aufzufinden. Am 8. Aug. hatte Hr. Prof.
\textsc{Harding} die Gegend in den Fischen durchmustert, wo er stehen musste, und am 9. Aug. erkannten wir ihn sofort in einem Sterne 9. Grösse, der sich merklich nach Süden bewegt hatte. Den 10. Aug. wurde diese Wiederauffindung zur völligen Gewissheit gebracht. Hier sind die Resultate meiner Beobachtungen von diesen Tagen.


\begin{center}\fontsize{9}{11}\selectfont\setlength{\tabcolsep}{5pt}\renewcommand{\arraystretch}{1.15}\begin{tabular}{lrrr} 1809 & \shortstack{Mittlere Zeit\\in Göttingen.} & \shortstack{Scheinb. gerade\\Aufsteig. der $\PallasSymbol$} & \shortstack{Scheinb. nördl.\\Abweichung.}\\\hline Aug. 9 & $11^{\text{u}}\;42^{\text{m}}\;48^{\text{s}}$ & $7^\circ\;51'\;44''$ & $2^\circ\;41'\;6''$\\ 10 & $11\;9\;58$ & $7\;50\;35$ & $2\;32\;37$\\
\end{tabular}\end{center}


\SourceLeaf{320}\par\indent Es erhellet hieraus, dass die Ephemeride im astron. Jahrbuche die Declinationen jetzt gut, aber die Rectascensionen um $20'$ zu klein gibt.

\par\indent Am 10. Aug. habe ich auch die Ceres aufgefunden; ein paar noch nicht vollständig reducirte Beobachtungen zeigen wenigstens, dass meine Ephemeride im Maiheft der M. C. den Ort auf die Minute genau darstellt. Die Ceres ist 8. bis 9. Grösse.

\par\indent In der Nacht vom 12. auf den 13. glaubte Hr. Prof. \textsc{Harding} auch bereits die Vesta als einen Stern 9. bis 10. Grösse sehr nahe auf dem Platze zu erkennen, welchen die vom Hrn. Dr.
\textsc{Schumacher} berechnete Ephemeride angibt; dies bedarf indess erst noch der Bestätigung, wozu die erste heitere Nacht benutzt werden wird.

\clearpage
\fancyhead[L]{\small PALLAS. CERES. VESTA.}
\begin{center}
\textbf{PALLAS. CERES. VESTA.}
\par\vspace{4pt}\hrule\vspace{4pt}
\fontsize{10}{12}\selectfont Göttingische gelehrte Anzeigen. Stück 136. Seite 1348..1350. 1809 August 26.
\par\vspace{4pt}\hrule\vspace{4pt}
\end{center}

\par\indent Von der in der Nacht vom 9. auf den 10. August auf der hiesigen Sternwarte geschehenen Wiederauffindung der \emph{Pallas} wird eine kurze Nachricht in unsern Blättern den Astronomen nicht unwillkommen sein, zumal da der Mangel an brauchbaren Beobachtungen aus dem vorigen Jahre den Hrn. Prof. \textsc{Gauss} abgehalten hatte, eine Verbesserung der Elemente vorzunehmen, und man sich daher zur Wiederauffindung des Planeten noch an die ältern Elemente halten muss, nach welchen Hr. Prof.
\textsc{Bode} im astronomischen Jahrbuche für 1811 eine kleine Ephemeride berechnet hat. Schon am 8. August wurde die Gegend, wo die Pallas erwartet wurde, durchmustert, und am folgenden Abend zeigte sich einer von den aufgezeichneten Sternen beträchtlich nach Süden vorgerückt. Am 10. August bestätigte es sich vollkommen, dass dies die Pallas gewesen sei. Folgende zwei Beobachtungen sind bisher vom Hrn. Prof. \textsc{Gauss} angestellt worden: [dieselben wie die an von \textsc{Lindenau} im Briefe vom 14. Aug. 1809 mitgetheilten S. 314 d. B.] Die erste Declination ist nicht ganz so zuverlässig, als die zweite; die Rectascensionen können beide für gut gelten. Hiernach gibt die erwähnte Ephemeride von Hrn. Prof. \textsc{Bode} die Declination gut, aber die Rectascension um 20 Min. zu klein. Die Pallas erscheint übrigens bereits vollkommen mit der Helligkeit eines Sterns neunter Grösse.

\par\indent Am 10. August wurde auch die \emph{Ceres} wieder aufgefunden, und vom Hrn. Prof. \textsc{Gauss} einige Mal mit nahe stehenden Sternen verglichen. Diese Beobachtungen sind aber noch nicht vollständig reducirt. Inzwischen ist es hinreichend, zu bemerken, dass dieser Planet bis auf die Minute genau auf dem Platze gefunden wurde, wo er nach den XIII. Elementen (s. das 87. St. dieser gel. Anz. d. J.) [1809 Juni 3. S. 313 d. B.] und der darnach berechneten Ephemeride (im Mai-Hefte der Monatl. Correspondenz) erwartet wurde. Die Ceres erscheint gleichfalls in dem Lichte eines Sterns neunter Grösse.

\par\indent Auch die \emph{Vesta} glaubte Hr. Prof. \textsc{Harding} bereits am 12. August in einem feinen Sternchen neunter Grösse zu erkennen, sehr nahe bei dem Platze, wo sie nach den IV. Elementen (s. das 55. St. dieser Anz.) [1809 April 8. S. 311 d. B.] und der darnach vom Hrn. Doctor \textsc{Schumacher} berechneten Ephemeride (Monatl. Corresp. am angef. O.) stehen soll. Jedoch hat seitdem die Witterung noch nicht erlaubt, dies näher zu bestätigen.
